\documentclass[preprint,12pt,authoryear]{elsarticle}

\usepackage{amsmath,amssymb,amsthm}
\usepackage{booktabs}
\usepackage{array}
\usepackage{tabularx}
\usepackage{graphicx}
\graphicspath{{./}{figures/}}
\usepackage{enumitem}
\usepackage[colorlinks=true,urlcolor=blue,linkcolor=blue,citecolor=blue]{hyperref}
\emergencystretch=3em
\setlength{\tolerance}{2000}

\newtheorem{proposition}{Proposition}
\newtheorem{corollary}{Corollary}
\newcommand{\pos}[1]{\left[#1\right]_+}
\newcommand{\dconv}{\Delta b_{\mathrm{conv}}}
\newcommand{\Ups}{\Upsilon_c}
\newcommand{\Eceil}{E_{\mathrm{ceil}}}
\newcommand{\Acmin}{A_c^{\min}}
\newcommand{\tconv}{\tau_{\mathrm{conv}}}

\journal{Ecological Modelling}

\begin{document}

\begin{frontmatter}

\title{A rising biocapacity total can mask ecological-capital land drawdown: a two-land model}

\author[1]{Amin Abaee\corref{cor1}}
\cortext[cor1]{Corresponding author. Email: \href{mailto:amin_abaee@ut.ac.ir}{amin\_abaee@ut.ac.ir}. ORCID: \href{https://orcid.org/0000-0002-0019-1842}{0000-0002-0019-1842}.}
\affiliation[1]{organization={Independent Researcher}, city={Tehran}, country={Iran}}

\begin{abstract}
Biocapacity is reported as an aggregate, despite wide differences in yield and regeneration among its constituent land classes. A rising total alone cannot reveal whether its composition has shifted. We examine this accounting problem with a two-land model that includes delayed ecological regeneration, population response, and an explicit projection of the quality state onto an imposed admissible interval. The quality ceiling is a model-imposed admissibility bound, not an estimated ecological threshold. Under the model's baseline conditions, a shortfall-triggered transfer to fast-provisioning land with a higher annual contribution per hectare adds to model biocapacity through the conversion term even as capital-land area falls. In a representative run, total model biocapacity rises $39\%$ as capital-land area falls $46.16\%$; the modeled capital-land contribution falls $51.72\%$. On the reduced fixed-demand comparison, the equality relation reaches the minimum-area floor without a fold or critical slowing; this result does not characterize the separate surplus no-conversion family or every trajectory of the delayed model. Separately, NFA Edition 2026 account values for 1961--2023 show total world biocapacity up $24.86\%$, while $B/P$ falls $52.71\%$ and the production footprint per person rises $16.73\%$. Cropland contributes $+24.06$ percentage points, Forest biocapacity $-5.95$, and built-up land $+7.44$ as an accounting class. These account movements are not measurements of physical stocks. Component values help interpret the aggregate.

\end{abstract}

\begin{keyword}
\footnotesize
biocapacity \sep ecological footprint \sep land-use conversion \sep aggregation \sep identifiability \sep critical slowing down \sep delay differential equations \sep carrying capacity
\end{keyword}

\end{frontmatter}

\section{Introduction}\label{sec:intro}

In ecology, \emph{carrying capacity} describes the population a resource base can sustain. Many models treat it as a fixed ceiling, in part because it is difficult to estimate \citep{Cohen1995}. If resource use changes the basis of future production, however, a fixed ceiling may not capture that change. This raises a measurement question: can a headline environmental metric report a rising total while a productive land class loses area?

The National Footprint and Biocapacity Accounts (NFA) report biological supply as biocapacity and human demand as the Ecological Footprint. Both are expressed in global hectares (gha), which adjust area for biological productivity relative to the world average in a given year; the annual basis is implicit \citep{GFNStandards2009,GFNGuidebook2019}. Because a global hectare combines area and productivity, a change in the total does not by itself show whether yields, area, or the mix of land uses changed. Has productivity increased, has the area in production shifted, or has the composition of the account changed?

In composite-indicator theory, this ambiguity is called \emph{compensation}: a gain in one component can offset a loss in another \citep{Nardo2008,Becker2017,Fischer2022}. The aggregate remains useful, but reading it alone can conceal component changes. Related critiques discuss the interpretation of the Ecological Footprint and biocapacity accounts \citep{Blomqvist2013,GiampietroSaltelli2014,vdBerghGrazi2015,Galli2016}.

Most human--environment models summarize nature as an aggregate stock or a unified feedback. Examples include the \emph{World3} simulations \citep{Meadows1972,Meadows1974,Meadows2004}, models of Easter Island \citep{BranderTaylor1998,Anderies2000}, and coupled Earth-system frameworks \citep{Motesharrei2016,Donges2017,CalvinBondLamberty2018,HendersonLoreau2018}. Other work debates whether human carrying capacity is fixed, variable, or useful as a concept \citep{Cohen1995,vdBerghRietveld2004,Simon1996,Meyer1994,MeyerAusubel1999,Dasgupta2021}; global demographic assessments often project population separately from environmental constraints \citep{UN2022}. Land-use studies, meanwhile, document conversion among forests, grasslands, wetlands, and cropland \citep{Haberl2007,Krausmann2013,Erb2017}. Together, these literatures motivate a narrower question: what can one biocapacity total reveal when land classes differ in yield and regeneration?

\paragraph{Analogies} To make the accounting issue concrete, imagine an orchard with fast-provisioning ground, such as intensively managed beds, beside slower-renewing ground with trees and living soil. Annual harvest is a \emph{flow}; trees, roots, and soil are assets that support future flows. In our model, a shortfall can trigger a transfer from ecological-capital land to fast-provisioning land. If the receiving class has a higher annual contribution per hectare in that scenario, the conversion term adds to $B$ at the margin while capital-land area $A_c$ falls; other flows can offset that increment. This is a conditional model mechanism, not evidence that NFA accounts measure physical stocks or that a particular real-world conversion occurred. Replanting may restore a later harvest, but it is not a one-for-one measure of the condition or recovery of the distinct assets represented by the model's capital-land class. The analogy separates current output from the state of its productive base; it does not assert that the classes map directly to NFA categories.

A second analogy separates observed service from structural condition: an elevator rated for ten people might continue operating while carrying fourteen, but the ride itself does not establish whether repeated overloading has changed the cable's safety margin. Continued service is not a measurement of structural condition, nor does the analogy predict inevitable failure. Related material-ledger work discusses a broader flow--stock masking problem; here we study its composition-specific form in a biocapacity model \citep{Abaee2026c}.

\paragraph{Research questions} This paper addresses three questions:
\begin{enumerate}
  \item \textbf{The masking dynamic:} Under what conditions can conversion raise total model biocapacity while $A_c$ declines, how long can this divergence last, and can the aggregate identify the land shift?
  \item \textbf{System trajectory:} Under the model's stated assumptions, does drawdown involve a fold, hysteresis, or critical slowing, or does the reduced comparison follow another path?
  \item \textbf{Limits and recovery:} What limits drawdown under fixed demand, and what modelled process allows capital-land area to recover?
\end{enumerate}

\paragraph{Core contributions} We make three contributions:
\begin{itemize}
  \item We construct a dynamic two-land model that conserves total area, represents capital-land quality separately from area, and includes delays in ecological regeneration and demographic response.
  \item We derive the condition under which a conversion term raises aggregate biocapacity while $A_c$ falls. In a reduced fixed-demand comparison, the balance relation reaches the area floor without a fold or critical slowing; realized paths in the delayed model remain history-dependent. These are conditional model results, not general claims about observed ecosystems.
  \item We decompose NFA Edition 2026 account values for 1961--2023 and examine a separate cropland-area/yield proxy. The account series show a compositional change, not a census of physical ecological stocks.
\end{itemize}

Finally, we derive five testable empirical predictions. The model parameters are representative rather than fitted, so simulations illustrate conditional mechanisms rather than historical forecasts. Aggregate biocapacity is therefore most informative when read alongside its component accounts.

